\documentclass[10pt,a4paper]{amsart}
\usepackage[margin=19mm]{geometry}
\usepackage{amsmath,amssymb,mathtools}
\usepackage[scr=rsfs,cal=euler]{mathalfa}
\usepackage{xfrac}
\usepackage[unicode,hidelinks,bookmarks=false]{hyperref}
\newcommand{\ie}{i.e.\ }
\newcommand{\cf}{cf.\ }
\newcommand{\resp}{respectively\ }
\newcommand{\Subsection}{\subsection}
\providecommand{\nobreakdash}{\nobreak\hbox{-}}
\setlength{\emergencystretch}{2em}
\multlinegap0pt
\usepackage{bm}
\usepackage{bbm}
\usepackage{dsfont}
\usepackage{xspace}
\usepackage{cancel}
\usepackage{mathtools}
\usepackage{cleveref}
\usepackage{amsthm}
\makeatletter
\@ifundefined{theorem}{\newtheorem{theorem}{Theorem}}{}
\@ifundefined{corollary}{\newtheorem{corollary}[theorem]{Corollary}}{}
\@ifundefined{proposition}{\newtheorem{proposition}[theorem]{Proposition}}{}
\makeatother
\newcommand{\eval}{\operatorname{eval}}
\title[Confluent Vandermonde matrix and related topics]{Confluent Vandermonde matrix and related topics}
\author{Chi-Kwong Li, Jephian C.-H. Lin}
\date{Source: arXiv:2403.01474v3 (CC BY 4.0)}
\begin{document}
\maketitle

\noindent\emph{Source and licence.} Confluent Vandermonde matrix and related topics. Version arXiv:2403.01474v3 (CC BY 4.0), distributed under the Creative Commons Attribution 4.0 International licence, as recorded in the arXiv licence metadata for this exact version and shown on the abstract page https://arxiv.org/abs/2403.01474v3. Full rights notice: licenses/arxiv-2403.01474-CC-BY-4.0.txt.\par\smallskip

\noindent\textit{Editorial scope.} Counted unit: the Proposition whose source label is \texttt{prop:prodrule} in the article (source0.tex lines 744--752, source label \texttt{prop:prodrule}), delivered together with the article's own complete original proof of it (source0.tex lines 753--770, one \texttt{proof} environment of 18 lines). No mathematics is written, repaired or re-derived: statement and proof are the article's own text line for line. Required context retained uncounted: cor:hasseshift (lines 355--366). Every definition, display and label of those lines is retained unchanged. Omitted from this package and not redistributed: 1--354, 367--743, 771--989. Nothing inside a retained excerpt is omitted or altered: every retained body is delivered exactly as the article's lines are written, so no omission receipt of any kind accompanies this package. Citations: the retained excerpts carry no citation command at all, so no bibliography entry of the article is transcribed and no reference is redistributed. Reference adaptation: none was needed and none was made; every reference inside the delivered unit resolves inside it, so no unresolved reference or citation is produced. Printed slips kept literal: no printed word, symbol, hypothesis or number of the counted statement or of its proof was altered. Layout and class: the article's own class is \texttt{article} with packages including amsmath, amssymb, amsthm, hyperref, cite, cleveref, tikz, soul, cancel; the wrapper is the lane's pinned \texttt{amsart} editorial wrapper at 10pt on A4, which restates the theorem-like environments the retained text needs and adds the article's own macro definitions that the retained text needs (\texttt{\textbackslash eval}), with extra wrapper packages (bm, bbm, dsfont, xspace, cancel, mathtools, cleveref). The delivered numbering is the unit's own. Assets: no retained span contains a figure environment, a graphics-inclusion command, a PDF-page inclusion, a table or a TikZ picture; no image file is copied into this package, the package declares no asset dependency and no image file is redistributed with it. Licence: version arXiv:2403.01474v3 of this article is distributed under the Creative Commons Attribution 4.0 International licence, as recorded in the arXiv licence metadata for this exact version and shown on the abstract page https://arxiv.org/abs/2403.01474v3. A full rights notice is delivered at licenses/arxiv-2403.01474-CC-BY-4.0.txt. Characters: every retained excerpt is pure ASCII, so the delivered engine, XeLaTeX with its default Latin Modern text font, reports no missing character.
\begin{corollary}
\label{cor:hasseshift}
Let $p(x)$ be a polynomial in $\mathbb{F}[x]$.  Then 
\[
    \mathcal{D}^{(r)}(x - \lambda)^m p(x) \big|_{x = \lambda} = 
    \begin{cases}
        \mathcal{D}^{(r - m)}p(\lambda) & \text{ if } r \geq m, \\
        0 & \text{ otherwise.}
    \end{cases}
\]

\end{corollary}

\begin{proposition}
\label{prop:prodrule}
Let $\Lambda$ be a multiset and $f(x)\in {\mathbb F}[x]$. Then
\begin{enumerate}
    \item $\mathcal{D}^{(0)}(xf(x))\big|_{x = \lambda} = \lambda \mathcal{D}^{(0)}(f(x))\big|_{x = \lambda}$, 
    \item $\mathcal{D}^{(r)}(xf(x))\big|_{x = \lambda} = \mathcal{D}^{(r-1)}(f(x))\big|_{x = \lambda} + \lambda \mathcal{D}^{(r)}(f(x))\big|_{x = \lambda}$ if $r \geq 1$, and   
    \item $\eval_\Lambda(xf(x)) = J_\Lambda \eval_\Lambda(f(x))$.
\end{enumerate}
\end{proposition}

\begin{proof}
The first identity holds since both sides equal to $\lambda f(\lambda)$.  

For the field of real numbers, the second identity could be written as 
\[
    \frac{1}{r!}\cdot \frac{d^r}{dx^r} (xf(x))\Big|_{x = \lambda} = \frac{1}{r!}\left( \binom{r}{1}\frac{d^{r-1}}{dx^{r-1}} (f(x)) + x\cdot \frac{d^r}{dx^r} (xf(x)) \right)\Big|_{x = \lambda}
\]
and is the product rule applied to $xf(x)$.  However, we emphasize that the identity remains valid from the Hasse derivative point of view, and the proof is even more intuitive.  By direct computation and \cref{cor:hasseshift}, we have 
\[
    \begin{aligned}
        \mathcal{D}^{(r)}(xf(x))\big|_{x = \lambda} &= \mathcal{D}^{(r)}((x - \lambda)f(x) + \lambda f(x))\big|_{x = \lambda} \\
        &= \mathcal{D}^{(r)}((x - \lambda)f(x))\big|_{x = \lambda} + \mathcal{D}^{(r)}(\lambda f(x))\big|_{x = \lambda} \\
        &= \mathcal{D}^{(r-1)}(f(x))\big|_{x = \lambda} + \lambda\mathcal{D}^{(r)}(f(x))\big|_{x = \lambda}.
    \end{aligned}
\]

Finally, the third identity is simply a matrix-vector form of the first two identities.  
\end{proof}

\end{document}
